\ChapterDivider
  {Network Theory}
  {Basic Concepts of Networks}
  {\begin{enumerate}[leftmargin=*,itemsep=2pt,topsep=2pt]
  \item Introduction to Electrical Quantities
    \item Basic Laws and Classification of Network Elements
    \item Resistors, Capacitors and Inductors
    \item Energy Sources in Electrical Circuits
    \item Series-Parallel connections
    \item Methods of Analyzing Electrical Circuits - Nodal and Mesh Analysis
  \end{enumerate}}

%----------Q1--------------%
\MCQ{1}{2026}{1}{1}{A}{
Refer to the four circuits shown.
\QuestionFigure[1.0\columnwidth]{img/ch1_basic_concepts_of_networks/Q1_26_1.png}

\noindent Which one of the following options for $k_1$, $k_2$, $k_3$, and $k_4$ makes all of them realizable?
\begin{choices}
  \Option{$k_1 = 1, k_4 = -1/3$, for all values of $k_2$ and $k_3$}
  \Option{$k_2 = -2, k_3 = +1/3$, for all values of $k_1$ and $k_4$}
  \Option{$k_1 = 2, k_2 = 0.5, k_3 = -2/3, k_4 = -3$}
  \Option{$k_1 = 2, k_2 = -0.5, k_3 = -2/3, k_4 = +3$}
\end{choices}
}

%----------Q2--------------%
\MCQ{1}{2026}{2}{1}{B}{
A circuit with ideal elements is shown.
\QuestionFigure[0.6\columnwidth]{img/ch1_basic_concepts_of_networks/Q1_26_2.png}

\noindent Which one of the following options correctly identifies all the linear elements in the circuit?
\InlineOptions
  {R only}
  {R, L, and C only}
  {D only}
  {L, C, and D only}
}


%----------Q3--------------%
\MCQ{1}{2025}{1}{1}{B}{
A nullator is defined as a circuit element where the voltage across the device and the current through the device are both zero. A series combination of a nullator and a resistor of value, $R$, will behave as a
\InlineOptions
  {resistor of value $R$}
  {nullator}
  {open circuit}
  {short circuit}
}

%----------Q4--------------%
\MCQ{1}{2025}{2}{1}{B}{
The I-V characteristics of the element between the nodes $X$ and $Y$ is best depicted by
\QuestionFigure[0.4\columnwidth]{img/ch1_basic_concepts_of_networks/Q1_25_2.png}
\InlineOptions
  {\QuestionFigureNoNumber[width=0.35\columnwidth]{img/ch1_basic_concepts_of_networks/Q1_25_2_1.png}}
  {\QuestionFigureNoNumber[width=0.35\columnwidth]{img/ch1_basic_concepts_of_networks/Q1_25_2_2.png}}
  {\QuestionFigureNoNumber[width=0.35\columnwidth]{img/ch1_basic_concepts_of_networks/Q1_25_2_3.png}}
  {\QuestionFigureNoNumber[width=0.35\columnwidth]{img/ch1_basic_concepts_of_networks/Q1_25_2_4.png}}
}
%----------Q5--------------%
\MCQ{1}{2024}{1}{1}{A}{
The number of junctions in the circuit is
\QuestionFigure[0.6\columnwidth]{img/ch1_basic_concepts_of_networks/Q1_24_1.png}
\InlineOptionsOneLine
  {6}
  {7}
  {8}
  {9}
}

%----------Q6--------------%
\MCQ{1}{2024}{2}{1}{A}{
All the elements in the circuit are ideal. The power delivered by the $10\text{ V}$ source in watts is
\newpage
\QuestionFigure[0.9\columnwidth]{img/ch1_basic_concepts_of_networks/Q1_24_2.png}
\InlineOptions
  {0}
  {50}
  {100}
  {dependent on the value of $\alpha$}
}
%----------Q7--------------%
\NAT{1}{2023}{1}{1}{5.9 to 6.1}{
For the circuit shown in the figure, $V_1 = 8\text{ V, DC}$ and $I_1 = 8\text{ A, DC}$. The voltage $V_{\text{ab}}$ in Volts is \underline{\hspace{1.5cm}}(Round off to 1 decimal place).
\QuestionFigure[0.9\columnwidth]{img/ch1_basic_concepts_of_networks/Q1_23_1.png}
}
%----------Q8--------------%
\NAT{1}{2022}{1}{2}{98 to 102 }{
In the circuit shown below, the magnitude of voltage $V_1$ in volts, across the $8\text{ k}\Omega$ resistor is \underline{\hspace{1.5cm}}. (round off to nearest integer)
\QuestionFigure[1.0\columnwidth]{img/ch1_basic_concepts_of_networks/Q1_22_1.png}
}
%----------Q9--------------%
\NAT{1}{2021}{1}{1}{-3.30 to -3.20}{
In the given circuit, for voltage $V_y$ to be zero, value of $\beta$ should be \underline{\hspace{1.5cm}}. (Round off to 2 decimal places).
\QuestionFigure[0.9\columnwidth]{img/ch1_basic_concepts_of_networks/Q1_21_1.png}
}
%----------Q10--------------%


\NAT{1}{2019}{1}{1}{1.3 to 1.5}{
The current $I$ flowing in the circuit shown below in amperes (round off to one decimal place) is \underline{\hspace{1.5cm}}.
\QuestionFigure[0.9\columnwidth]{img/ch1_basic_concepts_of_networks/Q1_19_1.png}
}


%----------Q13--------------%
\NAT{1}{2017}{1}{1}{248 to 252}{
The power supplied by the $25\text{ V}$ source in the figure shown below is \underline{\hspace{1.5cm}}\text{ W}.
\QuestionFigure[1.0\columnwidth]{img/ch1_basic_concepts_of_networks/Q1_17_1.png}
}
%----------Q14--------------%
\MCQ{1}{2016}{1}{1}{D}{
In the circuit shown below, the voltage and current sources are ideal. The voltage ($V_{\text{out}}$) across the current source, in volts, is
\QuestionFigure[0.9\columnwidth]{img/ch1_basic_concepts_of_networks/Q1_16_1.png}
\InlineOptionsOneLine
  {0}
  {5}
  {10}
  {20}
}

%----------Q15--------------%
\NAT{1}{2016}{2}{1}{0.5 to 0.5}{
In the given circuit, the current supplied by the battery, in ampere, is \underline{\hspace{1.5cm}}.
\QuestionFigure[0.9\columnwidth]{img/ch1_basic_concepts_of_networks/Q1_16_2.png}
}
%----------Q16--------------%
\NAT{1}{2016}{3}{1}{1.9 to 2.1}{
In the portion of a circuit shown, if the heat generated in $5\ \Omega$ resistance is $10\text{ calories per second}$, then heat generated by the $4\ \Omega$ resistance, in calories per second, is \underline{\hspace{1.5cm}}.
\QuestionFigure[0.8\columnwidth]{img/ch1_basic_concepts_of_networks/Q1_16_3.png}
}
%----------Q17--------------%
\MCQ{1}{2016}{4}{1}{D}{
The graph associated with an electrical network has 7 branches and 5 nodes. The number of independent KCL equations and the number of independent KVL equations, respectively, are
\InlineOptionsOneLine
  {2 and 5}
  {5 and 2}
  {3 and 4}
  {4 and 3}
}

%----------Q18--------------%
\NAT{1}{2015}{1}{2}{0.48 to 0.52}{
In the given circuit, the parameter $k$ is positive, and the power dissipated in the $2\ \Omega$ resistor is $12.5\text{W}$. The value of $k$ is \underline{\hspace{1.5cm}}.
\QuestionFigure[0.85\columnwidth]{img/ch1_basic_concepts_of_networks/Q1_15_1.png}
}

%----------Q19--------------%
\MCQ{1}{2015}{2}{1}{A}{
The voltages developed across the $3\ \Omega$ and $2\ \Omega$ resistors shown in the figure are $6\text{V}$ and $2\text{V}$ respectively, with the polarity as marked. What is the power (in Watt) delivered by the $5\text{V}$ voltage source?
\QuestionFigure[0.85\columnwidth]{img/ch1_basic_concepts_of_networks/Q1_15_2.png}
\InlineOptionsOneLine
  {5}
  {7}
  {10}
  {14}
}

%----------Q20--------------%
\NAT{1}{2015}{3}{1}{0}{
The current $i$ (in Ampere) in the $2\ \Omega$ resistor of the given network is \underline{\hspace{1.5cm}}.
\newpage
\QuestionFigure[0.9\columnwidth]{img/ch1_basic_concepts_of_networks/Q1_15_3.png}
}


%----------Q21--------------%
\NAT{1}{2014}{1}{2}{10 to 10}{
In the figure, the value of resistor $R$ is $(25 + I/2)$ ohms, where $I$ is the current in amperes. The current $I$ is \underline{\hspace{1.5cm}}.
\QuestionFigure[0.4\columnwidth]{img/ch1_basic_concepts_of_networks/Q1_14_1.png}
}
%----------Q22--------------%
\NAT{1}{2014}{2}{2}{2469 to 2473}{
An incandescent lamp is marked $40\text{W}$, $240\text{V}$. If resistance at room temperature ($26^{\circ}\text{C}$) is $120\ \Omega$ and temperature coefficient of resistance is $4.5 \times 10^{-3} / ^{\circ}\text{C}$, then its 'ON' state filament temperature is $^{\circ}\text{C}$ is approximately \underline{\hspace{1.5cm}}.
}
%----------Q23--------------%
\NAT{1}{2014}{3}{1}{330 to 330}{
The three circuit elements shown in the figure are part of an electric circuit. The total power absorbed by the three circuit elements in watt is \underline{\hspace{1.5cm}}.
\QuestionFigure[0.6\columnwidth]{img/ch1_basic_concepts_of_networks/Q1_14_3.png}
}
%----------Q24--------------%
\MCQ{1}{2010}{1}{2}{B}{
If the $12\ \Omega$ resistor draws a current of $1\text{ A}$ as shown in the figure, the value of resistance $R$ is
\QuestionFigure[0.8\columnwidth]{img/ch1_basic_concepts_of_networks/Q1_10_1.png}
\InlineOptionsOneLine
  {$4\ \Omega$}
  {$6\ \Omega$}
  {$8\ \Omega$}
  {$18\ \Omega$}

}
%----------Q25--------------%
\MCQ{1}{2010}{2}{1}{B}{
As shown in the figure given below, a $1\ \Omega$ resistance is connected across a source that has a load line $v + i = 100$. The current through the resistance is
\QuestionFigure[0.7\columnwidth]{img/ch1_basic_concepts_of_networks/Q1_10_2.png}
\InlineOptionsOneLine
  {$25\text{ A}$}
  {$50\text{ A}$}
  {$100\text{ A}$}
  {$200\text{ A}$}

}
%----------Q26--------------%
\MCQ{1}{2009}{1}{2}{B}{
For the circuit shown below, find out the current flowing through the $2\ \Omega$ resistance. Also, identify the changes to be made to double the current through the $2\ \Omega$ resistance.
\QuestionFigure[0.8\columnwidth]{img/ch1_basic_concepts_of_networks/Q1_09_1.png}
\InlineOptions
  {($5\text{ A}$; Put $V_{\text{s}} = 20\text{ V}$)}
  {($2\text{ A}$; Put $V_{\text{s}} = 8\text{ V}$)}
  {($5\text{ A}$; Put $I_{\text{s}} = 10\text{ A}$)}
  {($7\text{ A}$; Put $I_{\text{s}} = 12\text{ A}$)}

}
%----------Q27--------------%
\MCQ{1}{2008}{1}{2}{B}{
In the circuit shown in the figure. The value of the current $I$ will be given by
\QuestionFigure[0.95\columnwidth]{img/ch1_basic_concepts_of_networks/Q1_08_1.png}
\InlineOptionsOneLine
  {$0.31\text{ A}$}
  {$1.25\text{ A}$}
  {$1.75\text{ A}$}
  {$2.5\text{ A}$}

}
%----------Q28--------------%
\MCQ{1}{2008}{2}{2}{A}{
Assuming ideal elements in the circuit shown below, the voltage $V_{\text{ab}}$ will be
\QuestionFigure[0.65\columnwidth]{img/ch1_basic_concepts_of_networks/Q1_08_2.png}
\InlineOptionsOneLine
  {$-3\text{V}$}
  {$0\text{V}$}
  {$3\text{V}$}
  {$5\text{V}$}

}

%----------Q30--------------%
\MCQ{1}{2007}{1}{2}{A}{
A $3\text{ V}$ dc supply with an internal resistance of $2\ \Omega$ supplies a passive non-linear resistance characterized by the relation $V_{\text{NL}} = I_{\text{NL}}^2$. The power dissipated in the non linear resistance is
\InlineOptionsOneLine
  {$1.0\text{ W}$}
  {$1.5\text{ W}$}
  {$2.5\text{ W}$}
  {$3.0\text{ W}$}
}


%----------Q31--------------%
\MCQ{1}{2005}{1}{1}{C}{
In the given figure, the value of $R$ is
\QuestionFigure[0.65\columnwidth]{img/ch1_basic_concepts_of_networks/Q1_05_1.png}
\InlineOptionsOneLine
  {$2.5\ \Omega$}
  {$5.0\ \Omega$}
  {$7.5\ \Omega$}
  {$10.0\ \Omega$}
}

%----------Q32--------------%
\MCQ{1}{2004}{1}{2}{C}{
In the given figure, the value of the source voltage $E$ is
\QuestionFigure[0.8\columnwidth]{img/ch1_basic_concepts_of_networks/Q1_04_1.png}
\InlineOptionsOneLine
  {$12\text{ V}$}
  {$24\text{ V}$}
  {$30\text{ V}$}
  {$44\text{ V}$}

}

%----------Q33--------------%
\MCQ{1}{2003}{1}{1}{C}{
Figure shows the waveform of the current passing through an inductor of resistance $1\ \Omega$ and inductance $2\text{ H}$. The energy absorbed by the inductor in the first four seconds is
\QuestionFigure[0.65\columnwidth]{img/ch1_basic_concepts_of_networks/Q1_03_1.png}
\InlineOptionsOneLine
  {$144\text{ J}$}
  {$98\text{ J}$}
  {$132\text{ J}$}
  {$168\text{ J}$}

}



%----------Q35--------------%
\MCQ{1}{2001}{1}{2}{B}{
Consider the star network shown in Figure. The resistance between terminals A and B with C open is $6\ \Omega$, between terminals B and C with A open is $11\ \Omega$, and between terminals C and A with B open is $9\ \Omega$. Then
\QuestionFigure[0.65\columnwidth]{img/ch1_basic_concepts_of_networks/Q1_01_1.png}
\begin{choices}
  \Option{$R_{\text{A}} = 4\ \Omega$, $R_{\text{B}} = 2\ \Omega$, $R_{\text{C}} = 5\ \Omega$}
  \Option{$R_{\text{A}} = 2\ \Omega$, $R_{\text{B}} = 4\ \Omega$, $R_{\text{C}} = 7\ \Omega$}
  \Option{$R_{\text{A}} = 3\ \Omega$, $R_{\text{B}} = 3\ \Omega$, $R_{\text{C}} = 4\ \Omega$}
  \Option{$R_{\text{A}} = 5\ \Omega$, $R_{\text{B}} = 1\ \Omega$, $R_{\text{C}} = 10\ \Omega$}
\end{choices}
}



%----------Q37--------------%
\NAT{1}{1999}{1}{2}{*}{
Solve the circuit shown in figure using the mesh method of analysis and determine the mesh currents $I_1$, $I_2$ and $I_3$. Evaluate the power developed in the $10\text{V}$ voltage source.
\QuestionFigure[0.8\columnwidth]{img/ch1_basic_concepts_of_networks/Q1_99_1.png}
}

%----------Q38--------------%
\MCQ{1}{1999}{2}{2}{B}{
When a resistor $R$ is connected to a current source, it consumes a power of $18\text{W}$. When the same $R$ is connected to a voltage source having the same magnitude as the current source, the power absorbed by $R$ is $4.5\text{W}$. The magnitude of the current source and the value of $R$ are
\InlineOptions
  {$\sqrt{18}\text{A}$ and $1\ \Omega$}
  {$3\text{A}$ and $2\ \Omega$}
  {$1\text{A}$ and $18\ \Omega$}
  {$6\text{A}$ and $0.5\ \Omega$}

}

%----------Q39--------------%
\MCQ{1}{1999}{3}{2}{C}{
The color code of a $1\text{k}\Omega$ resistance is:
\InlineOptions
  {black, brown, red}
  {red, brown, brown}
  {brown, black, red}
  {black, black, red}

}

%----------Q40--------------%
\NAT{1}{1997}{1}{2}{Inductor, 2H}{
The voltage and current waveforms for an element are shown in figure. The circuit element is \underline{\hspace{1.5cm}} And its value is \underline{\hspace{1.5cm}}.

\QuestionFigure[0.9\columnwidth]{img/ch1_basic_concepts_of_networks/Q1_97_1.png}
}
%----------Q41--------------%
\NAT{1}{1997}{2}{2}{31V, 13A}{
The values of $E$ and $I$ for the circuit shown in figure, are\dots\dots$\text{V}$ and \dots\dots.
\QuestionFigure[0.8\columnwidth]{img/ch1_basic_concepts_of_networks/Q1_97_2.png}
}
%----------Q42--------------%
\NAT{1}{1997}{3}{2}{5}{
A $10\text{V}$ battery with an internal resistance of $1\ \Omega$ is connected across a nonlinear load whose V-I characteristic is given by $7I = V^2 + 2V$. The current delivered by the battery is \underline{\hspace{1.5cm}}\text{A}.
}
%----------Q43--------------%
\MCQ{1}{1997}{4}{1}{C}{
Energy stored in a capacitor over a cycle, when excited by an a.c. source is
\begin{choices}
  \Option{The same as that due to a d.c. source of equivalent magnitude}
  \Option{Half of that due to a d.c. source of equivalent magnitude}
  \Option{Zero}
  \Option{None of the above}
\end{choices}
}
%----------Q44--------------%
\MCQ{1}{1997}{5}{1}{B}{
A practical current source is usually represented by
\begin{choices}
  \Option{Resistance in series with an ideal current source}
  \Option{A resistance in parallel with an ideal current source}
  \Option{A resistance in parallel with an ideal voltage source}
  \Option{None of the above}
\end{choices}
}
%----------Q45--------------%
\MCQ{1}{1997}{6}{1}{C}{
An ideal voltage source will charge an ideal capacitor
\InlineOptions
{In infinite time}
  {Exponentially}
  {Instantaneously}
  {None of these}
}
%----------Q46--------------%
\MCQ{1}{1996}{1}{1}{C}{
In the circuit shown in figure X is an element which always absorbs power. During a particular operations, it sets up a current of $1\text{ amp}$ in the direction shown and absorbs a power $P_{\text{x}}$. It is possible that X can absorb the same power $P_{\text{x}}$ for another current $i$, the value of this current is
\QuestionFigure[0.5\columnwidth]{img/ch1_basic_concepts_of_networks/Q1_96_1.png}
\InlineOptions
  {$(3-\sqrt{14})\text{ amps}$}
  {$(3+\sqrt{14})\text{ amps}$}
  {$5\text{ amps}$}
  {None of these}
}
%----------Q47--------------%
\MCQ{1}{1992}{1}{1}{C}{
All resistances in the circuit in figure are of $R$ ohms each. The switch is initially open. What happens to the lamp's intensity when the switch is closed?
\QuestionFigure[0.6\columnwidth]{img/ch1_basic_concepts_of_networks/Q1_92_1.png}
\InlineOptions
  {Increases}
  {Decreases}
  {Remains same}
  {Answer depends on the value of $R$}

}

\EndChapter
